% ============================================================
% APÊNDICE A — MATRIZ DE LITERATURA CONSOLIDADA
% ============================================================
\chapter*{Apêndice A — Matriz de Literatura Consolidada (14 Capítulos)}
\addcontentsline{toc}{chapter}{Apêndice A — Matriz de Literatura Consolidada}

Este apêndice apresenta a matriz global de referências utilizada na fundamentação teórica dos 14 capítulos. Cada capítulo possui 4–6 eixos temáticos com referências seminais (clássicos consolidados) e referências 2020–2025 (relatórios institucionais, artigos de alto impacto, normas atualizadas). A matriz completa em formato tabular está disponível no repositório do livro (\texttt{02-matriz-literatura/}).

\section*{Estrutura da Matriz por Capítulo}

\begin{longtable}{p{2.5cm}p{5.5cm}p{5.5cm}}
\toprule
\textbf{Capítulo} & \textbf{Eixos Temáticos} & \textbf{Referências-Chave (Seleção)} \\
\midrule
\endhead
\midrule
\multicolumn{3}{r}{\textit{Continua na próxima página}} \\
\endfoot
\bottomrule
\endlastfoot

1. Planejamento Estratégico &
Escolas de estratégia; Posicionamento (Porter); RBV \& Cap. Dinâmicas; Análise ambiente/cenários; Execução (BSC, OKR, Rolling Forecast); VUCA/BANI/Blue Ocean &
Chandler 1962; Ansoff 1965; Mintzberg 1998; Porter 1980/1985/1996; Barney 1991; Prahalad \& Hamel 1990; Teece et al. 1997; Teece 2007; Kaplan \& Norton 1992/1996/2001/2008; Grove 1983; Doerr 2018; Hope \& Fraser 2003; Schwartz 1991; Kim \& Mauborgne 2005/2015; Cascio 2020; Gartner 2023–2024; McKinsey 2022–2024; BCG 2024 \\

2. Gestão Financeira &
Teoria finanças corporativas; Análise demonstrações; Orçamento/Beyond Budgeting/Rolling Forecast; Fluxo caixa/capital giro; Valuation/EVA; Finanças digitais/Open Finance &
Modigliani \& Miller 1958/1963; Myers \& Majluf 1984; Damodaran 2020/2024; Koller et al. 2024; Rappaport 1986; Stewart 1991; Hope \& Fraser 2003; Jensen 2001; BCB Relatórios 2024; Febraban 2024; KPMG/Ibmec WACC 2024; World Bank Findex 2024; BCB Open Finance 2024 \\

3. Operações \& Cadeia Suprimentos &
Fundamentos operações; Lean/TPS; Seis Sigma; Cadeia suprimentos (design/resiliência/digital); ERP/APS/Operações 4.0; ESG Supply Chain &
Slack et al.; Chase et al.; Heizer \& Render; Ohno 1988; Womack \& Jones 1996; Liker 2004; Spear \& Bowen 1999; Chopra \& Meindl; Christopher; Fisher 1997; Lee 2004; Davenport 1998; Davenport \& Ronanki 2018; Gartner CSC 2024; WEF Resilient SC 2023; McKinsey SC Resilience 2024; CDP SC 2024; EcoVadis 2024 \\

4. Gestão Pessoas \& Liderança &
Comportamento organizacional/cultura; SHRM; Avaliação 360°/feedback contínuo; Talentos/sucessão/pipeline; Competências/carreiras; DEIP &
Robbins \& Judge; Schein 1985/2010; Beer et al. 1984; Guest 1987; Wright \& McMahan 1992; Ulrich 1997/2016; Aguinis; London \& Smither 1995; Rothwell; Charan et al. 2001; Groysberg; DDI 2023/2025; Dutra; Le Boterf; Fleury \& Fleury 2004; WEF FoJ 2023/2025; Deloitte HC Trends 2024; Gallup 2024; McKinsey 2022/2024; McKinsey Diversity 2023 \\

5. Marketing \& Vendas &
Fundamentos/holístico; Comportamento consumidor; STP; Mix 7Ps/serviços; Pesquisa mercado; CRM/Funil/RevOps; Marketing digital/MarTech/Growth &
Kotler \& Keller; Drucker; Levitt 1960; Kotler et al. 2021/2024; Kahneman 2011; McCarthy 1964; Ries \& Trout 1981; Booms \& Bitner 1981; Parasuraman et al. 1988; Payne \& Frow 2005; ChiefMarTec 2024; Gartner MarTech 2024; IAB AdInvest 2024; Salesforce/HubSpot/Clari/Gong 2024; Reforge Growth Loops \\

6. Gestão Projetos \& Portfólio &
PMBOK 7ª ed.; Ágil/Híbrido (Scrum/Kanban/SAFe); PMO/OPM3; PPM; Riscos projetos/entrega híbrida &
PMI PMBOK 7ª 2021; Kerzner; Meredith \& Mantel; Beck et al. 2001; Sutherland 2014; Schwaber \& Sutherland 2020; Anderson 2010; Leffingwell SAFe 2023; Hobbs \& Aubry 2007; Cooper et al.; Morris \& Pinto; PMI OPM3; PMI Pulse 2023/2024; Digital.ai Agile 2024; Gartner PPM 2024 \\

7. Inovação \& Desenvolvimento Produtos &
Teorias inovação; Disrupção/JTBD; Stage-Gate/Agile-Stage-Gate; Design Thinking/Service Design; Lean Startup/MVP; Ciclo vida/portfolio &
Schumpeter 1934/1942; Drucker 1985; Rogers 1962; Christensen 1997/2003/2016; Cooper 2011; Cooper \& Sommer 2016; Brown 2009; Liedtka \& Ogilvie 2011; Ries 2011; Blank 2013; Thomke 2020; Levitt 1965; Vernon 1966; Baghai et al. 1999; OECD Oslo Manual 2018; WIPO GII 2024; PINTEC 2024 \\

8. Gestão Qualidade \& Excelência &
Fundadores (Deming/Juran/Crosby/Feigenbaum); TQM/Kaizen; ISO 9001/MEG/Baldrige; Ferramentas 7+7/PDCA; Seis Sigma DMAIC; Qualidade 4.0 &
Deming 1986; Juran 1988; Crosby 1979; Feigenbaum 1991; Oakland; Imai 1986; Dahlgaard; Camp 1989; ISO 9001:2015; FNQ MEG 23ª 2023; NIST Baldrige 2023/2024; Ishikawa; Taguchi; ASQ Quality 4.0 2022; ISO 10012; FNQ Anuário 2024 \\

9. TI \& Dados &
SIG/MIS; Estratégia digital/TI vantagem; Governança TI/COBIT; Big Data/Analytics/Data-Driven; Cloud/DevOps; Segurança/LGPD &
Laudon \& Laudon; O'Brien; Porter \& Millar 1985; McFarlan \& McKenney 1983; Westerman et al. 2014; Rogers 2016; Bharadwaj et al. 2013; Iansiti \& Lakhani 2020; Weill \& Ross 2004; COBIT 2019; Davenport 2006/2018; Provost \& Fawcett 2013; Armbrust et al. 2010; Kim et al. 2016; ISO 27001:2022; NIST CSF 2.0 2024; LGPD/ANPD 2024; DORA 2024; Gartner 2024; IDC 2024; CGI.br 2024 \\

10. Comunicação, Stakeholders &
Comunicação organizacional; Stakeholder theory/saliência; Comunicação interna/engajamento; Mídias sociais/advocacy; Crise/reputação/sensemaking &
Shannon \& Weaver 1949; Argenti; Kunsch 2017; Grunig \& Hunt 1984; Freeman 1984; Mitchell et al. 1997; Frooman 1999; Clampitt 2016; Welch \& Jackson 2007; Kaplan \& Haenlein 2010; Kietzmann et al. 2011; Fombrun 1996; Coombs 2007; Weick et al. 2005; Edelman Trust 2024/2025; DataReportal 2024; Aberje 2024; Hootsuite/Sprout 2024; PwC Crisis 2024; RepTrak 2024 \\

11. Sustentabilidade, ESG &
RSE/TBL/CSV; ESG métricas/ratings/investimento; Relato integrado (GRI/ISSB/CSRD/TCFD); Economia circular/Net Zero; ODS/Agenda 2030; Clima/finanças verdes &
Carroll 1979/1991; Elkington 1997; Brundtland 1987; Porter \& Kramer 2011; PRI 2006; MSCI/Sustainalytics; SASB; Friede et al. 2015; IIRC 2021; GRI 2021; TCFD 2017; IFRS S1/S2 2023; EFRAG ESRS 2023/2024; SEC Climate 2024; CVM Res. 193 2023; EMF 2013/2024; McDonough \& Braungart 2002; SBTi 2023/2024; IPCC AR6 2023; NGFS 2024; ICMA 2024; B3 ISE 2024; BNDES Fundo Clima 2024 \\

12. Compliance \& Governança &
Teoria governança (agency/stewardship); Governança Brasil (Lei 6.404/Novo Mercado/IBGC); Compliance programas (ISO 37301/37001/DOJ); Anticorrupção (Lei 12.846/FCPA/UKBA); Controles internos/Three Lines; Ética/canais denúncia &
Jensen \& Meckling 1976; Donaldson \& Davis 1991; Shleifer \& Vishny 1997; Cadbury 1992; King IV 2016; Lei 6.404/1976; B3 Novo Mercado; IBGC Código 2024; OECD 2015; DOJ ECCP 2017/2023/2024; ISO 37301:2021; ISO 37001:2016; Lei 12.846/2013; Decreto 11.129/2022; FCPA; UK Bribery Act 2010; COSO 2013/2017; IIA Three Lines 2020; Crane \& Matten; Treviño \& Nelson; EU Whistleblower Directive 2021; NAVEX 2024; ECI 2024; PwC Econ Crime 2024; CGU 2024 \\

13. Gestão Riscos Corporativos (ERM) &
Fundamentos risco/incerteza; ERM/apetite/integração estratégica; Riscos financeiros/operacionais/crédito; Riscos estratégicos/emergentes/resiliência; Cultura risco/Three Lines; Risco climático/ESG/transição &
Knight 1921; Bernstein 1996; March \& Shapira 1987; COSO ERM 2004/2017; ISO 31000:2018; Fraser \& Simkins 2021; Beasley et al. 2005; Kaplan \& Mikes 2012; Hull; Basel III; Jorion; Taleb 2007/2012; Slywotzky \& Drzik 2005; WEF Resilience; ISO 22316 2017/2024; IIA Three Lines 2020; IRM Risk Culture 2022; TCFD 2017; NGFS 2023/2024; BCB CMN 4.947; ISSB S2 2023; CVM Res. 193; BCB Taxonomia Verde 2024; WEF Global Risks 2025; AON 2023; Protiviti 2024 \\

14. Transformação Digital \& IA &
Estratégia digital/transformação; IA preditiva/generativa; GenAI empresas; Plataformas/ecossistemas; Automação/hiperautomação/futuro trabalho; Marco legal/ética/IA responsável Brasil &
Westerman et al. 2014; Rogers 2016; Bharadwaj et al. 2013; Iansiti \& Lakhani 2020; Agrawal et al. 2018; Davenport \& Ronanki 2018; Brynjolfsson \& McAfee 2014; Russell \& Norvig 2020; Bommasani et al. 2021; Kaplan et al. 2020; McKinsey GenAI 2023/2024; Gartner GenAI 2024; NIST AI RMF 2023; OECD AI 2019/2024; EU AI Act 2024; Parker et al. 2016; Chui et al. 2016; Acemoglu \& Restrepo 2018; WEF FoJ 2023/2025; Gartner Hyperautomation 2024; LGPD 2018; Marco Civil 2014; PLC 2338/2023; CGI.br Diretrizes IA 2023; MCTI EBIA 2021/2024; ANPD Guia IA 2024; TCU Auditoria Algoritmos 2024; Stanford AI Index 2024/2025; CGI.br TIC Empresas 2024 \\

\end{longtable}

\section*{Critérios de Seleção das Referências 2020–2025}
\begin{itemize}
    \item \textbf{Relatórios institucionais:} McKinsey Global Institute, BCG Henderson Institute, Deloitte Insights, PwC, KPMG, Gartner, Forrester, IDC, WEF, OECD, IBGC, FGV, CNI, IBGE, BCB, CVM, B3, CGI.br, ANPD, ISO, COSO, IIA, PRI, GRI, ISSB, SBTi, TCFD, NGFS, NIST, AON, Marsh, Protiviti, S\&P, Fitch, Moody's.
    \item \textbf{Artigos acadêmicos alto impacto:} Strategic Management Journal, Academy of Management Review/Journal, Harvard Business Review, MIT Sloan Management Review, Journal of Finance, Journal of Financial Economics, MIS Quarterly, Information Systems Research, Journal of Marketing, Journal of Operations Management, Production and Operations Management, Risk Analysis, Journal of Risk and Financial Management.
    \item \textbf{Normas e regulamentos:} ISO (9001, 14001, 27001, 31000, 37001, 37301, 22316), COSO (IC, ERM), IIA (Three Lines, Standards), IFRS (S1, S2), GRI (Standards 2021), EFRAG (ESRS), TCFD, SASB, SBTi, NGFS, NIST (CSF 2.0, AI RMF), OECD (AI Principles, Compliance Guidance), EU (AI Act, CSRD, Whistleblower Directive), Brasil (LGPD, Lei 12.846, Decreto 11.129, Lei 13.964, CVM Res. 193, BCB Resoluções, PLC 2338/2023).
    \item \textbf{Dados estatísticos oficiais:} IBGE (PNAD, PINTEC, Contas Nacionais, ODS, Contas Ambientais), BCB (Focus, Rel. Estabilidade Financeira, Open Finance, Rel. Riscos Climáticos), CVM, B3 (ISE, Novo Mercado, Anuários), CNI (Sondagem, Panorama), Febraban (Panorama Bancário), ANBIMA, ILOS, ABILOG/ABRALOG, CGI.br (TIC Empresas/Domicílios), ANPD, CERT.br, SEEG, MCTI, BNDES, FINEP, TCU, CGU.
\end{itemize}

\section*{Observação Metodológica}
As referências listadas nos capítulos (seção 6 de cada capítulo) seguem a norma ABNT NBR 6023. DOIs e números de página foram omitidos intencionalmente na versão manuscrita para fluidez de leitura; serão inseridos na versão final para publicação. A matriz completa em Excel/CSV com 1.200+ entradas (capítulo, eixo, referência, tipo, ano, DOI, link) está disponível no repositório do projeto.
\end{flushright}